% QCM
\begin{enumerate}
    \item La molécule de chlorure d'hydrogène est polarisée selon~: \par
          \begin{itemize*}[label={\large$\square$},itemjoin={\hspace{2cm}}]
              \item \chemfig{%
                    \charge{90[anchor=180+\chargeangle]=$\scriptstyle\delta^{-}$}{H}-
                    \charge{90[anchor=180+\chargeangle]=$\scriptstyle\delta^{+}$}{Cl}}~;
              \item \chemfig{%
                    \charge{90[anchor=180+\chargeangle]=$\scriptstyle\delta^{+}$}{H}-
                    \charge{90[anchor=180+\chargeangle]=$\scriptstyle\delta^{-}$}{Cl}}.
          \end{itemize*}
    \item Le césium est un alcalin~; la formule du chlorure de césium (solide
          ionique) s'écrit~: \par
          \begin{itemize*}[label={\large$\square$},itemjoin={\hspace{2cm}}]
              \item \ce{CsCl}~;
              \item \ce{Cs2Cl}~;
              \item \ce{CsCl2}.
          \end{itemize*}
    \item Dans une solution de chlorure de calcium \ce{CaCl2} de concentration
          molaire $\qty{0.050}{\mol\per\L}$, la concentration
          molaire des ions calcium $\left[\ce{Ca^{2+}}\right]$ vaut~: \par
          \begin{itemize*}[label={\large$\square$},itemjoin={\hspace{2cm}}]
              \item $\qty{0.025}{\mol\per\L}$~;
              \item $\qty{0.050}{\mol\per\L}$~;
              \item $\qty{0.100}{\mol\per\L}$.
          \end{itemize*}
\end{enumerate}


